% ============================================================================
% ML INFERENCE PIPELINE
% ============================================================================
\section{ML Inference Pipeline}

The STAR robot uses machine learning for person detection as part of its safety system. This section documents the model selection and inference pipeline.

\subsection{Model Selection}

\begin{table}[H]
\centering
\caption{Person Detection Model Comparison}
\begin{tabular}{|l|c|c|c|c|}
\hline
\textbf{Model} & \textbf{Size (MB)} & \textbf{mAP@0.5} & \textbf{RPi5 FPS} & \textbf{Quantization} \\
\hline
YOLOv8n INT8 & 3.2 & 0.37 & 15-18 & INT8 \\
\hline
YOLOv5n FP16 & 7.5 & 0.35 & 8-10 & FP16 \\
\hline
MobileNet-SSD & 4.8 & 0.31 & 12-15 & INT8 \\
\hline
EfficientDet-Lite0 & 4.4 & 0.33 & 10-12 & INT8 \\
\hline
\end{tabular}
\end{table}

\textbf{YOLOv8n INT8} was selected for its optimal balance of accuracy and performance on the Raspberry Pi 5.

\subsection{Model Architecture}

\begin{itemize}
    \item \textbf{Base Model}: YOLOv8n (nano variant)
    \item \textbf{Input Resolution}: 320x320 (reduced from 640x640 for speed)
    \item \textbf{Quantization}: INT8 post-training quantization via TensorFlow Lite
    \item \textbf{Classes}: Single class (person) for focused detection
    \item \textbf{Inference Backend}: TensorFlow Lite with XNNPACK delegate
\end{itemize}

\subsection{Inference Pipeline}

\begin{lstlisting}[language=Python, caption=person\_detector.py]
import cv2
import numpy as np
from tflite_runtime.interpreter import Interpreter

class PersonDetector:
    def __init__(self, model_path: str = "yolov8n_int8.tflite"):
        # Load TFLite model with XNNPACK delegate for ARM optimization
        self.interpreter = Interpreter(
            model_path=model_path,
            num_threads=4  # Use all 4 Cortex-A76 cores
        )
        self.interpreter.allocate_tensors()

        self.input_details = self.interpreter.get_input_details()
        self.output_details = self.interpreter.get_output_details()
        self.input_shape = self.input_details[0]['shape'][1:3]  # (320, 320)

        # Detection thresholds
        self.confidence_threshold = 0.5
        self.nms_threshold = 0.4

    def preprocess(self, frame: np.ndarray) -> np.ndarray:
        """Resize and normalize input frame."""
        resized = cv2.resize(frame, tuple(self.input_shape))
        # INT8 quantization: scale to 0-255 range
        normalized = resized.astype(np.uint8)
        return np.expand_dims(normalized, axis=0)

    def detect(self, frame: np.ndarray) -> list:
        """Run inference and return person detections."""
        input_data = self.preprocess(frame)
        self.interpreter.set_tensor(
            self.input_details[0]['index'],
            input_data
        )
        self.interpreter.invoke()

        # Parse YOLO output format
        output = self.interpreter.get_tensor(
            self.output_details[0]['index']
        )
        detections = self._parse_output(output, frame.shape)
        return detections

    def _parse_output(self, output: np.ndarray,
                      original_shape: tuple) -> list:
        """Convert YOLO output to bounding boxes."""
        detections = []
        h, w = original_shape[:2]
        scale_x, scale_y = w / self.input_shape[1], h / self.input_shape[0]

        for detection in output[0]:
            confidence = detection[4]
            if confidence < self.confidence_threshold:
                continue

            # Extract bounding box (center format to corner format)
            cx, cy, bw, bh = detection[:4]
            x1 = int((cx - bw / 2) * scale_x)
            y1 = int((cy - bh / 2) * scale_y)
            x2 = int((cx + bw / 2) * scale_x)
            y2 = int((cy + bh / 2) * scale_y)

            detections.append({
                'bbox': (x1, y1, x2, y2),
                'confidence': float(confidence),
                'class': 'person'
            })

        # Apply non-maximum suppression
        return self._nms(detections)

    def _nms(self, detections: list) -> list:
        """Non-maximum suppression to remove overlapping boxes."""
        if not detections:
            return []

        boxes = np.array([d['bbox'] for d in detections])
        scores = np.array([d['confidence'] for d in detections])

        indices = cv2.dnn.NMSBoxes(
            boxes.tolist(),
            scores.tolist(),
            self.confidence_threshold,
            self.nms_threshold
        )

        return [detections[i] for i in indices.flatten()]
\end{lstlisting}

\subsection{Performance Optimization}

\begin{enumerate}
    \item \textbf{XNNPACK Delegate} -- ARM-optimized SIMD operations for 2-3x speedup over default TFLite.

    \item \textbf{Multi-threading} -- Utilize all 4 Cortex-A76 cores via \texttt{num\_threads=4}.

    \item \textbf{Reduced Resolution} -- 320x320 input (vs 640x640) reduces computation by 4x with minimal accuracy loss for person detection.

    \item \textbf{Frame Skipping} -- Process every 2nd frame (15 FPS effective) to maintain system responsiveness.

    \item \textbf{ROI Processing} -- In localization mode, only process regions where persons were previously detected.
\end{enumerate}

\subsection{Integration with Safety System}

The person detector feeds into the CV Safety System (Section~\ref{sec:cv-safety}):

\begin{lstlisting}[language=Python, caption=Integration with safety pipeline]
# In main control loop
detector = PersonDetector()
camera = cv2.VideoCapture(0)

while True:
    ret, frame = camera.read()
    if not ret:
        continue

    detections = detector.detect(frame)

    for det in detections:
        # Calculate distance from bounding box size
        # (calibrated for known person height)
        bbox_height = det['bbox'][3] - det['bbox'][1]
        estimated_distance = estimate_distance(bbox_height)

        # Trigger safety response based on distance
        if estimated_distance < EMERGENCY_STOP_DISTANCE:
            robot.emergency_stop()
        elif estimated_distance < SLOW_ZONE_DISTANCE:
            robot.reduce_speed(0.5)  # 50% max speed
\end{lstlisting}
